\documentclass[a4paper]{article}

% Basic packages
\usepackage[fontset=fandol]{ctex}
\usepackage{amsmath, amssymb}
\usepackage{graphicx}
\usepackage[margin=2.5cm]{geometry}
\usepackage[most]{tcolorbox}
\usepackage{etoolbox}
\usepackage{listings}
\usepackage{hyperref}
\usepackage{booktabs}
\usepackage{subcaption}
\usepackage{float}
\usepackage{tikz}
\IfFileExists{pgfplots.sty}{
    \usepackage{pgfplots}
    \pgfplotsset{compat=1.18}
}{}

% Highlight boxes
\newtcolorbox{knowledgebox}[1]{
    enhanced, colback=blue!5!white, colframe=blue!75!black, colbacktitle=blue!75!black,
    coltitle=white, fonttitle=\bfseries, title=#1, attach boxed title to top left={yshift=-2mm, xshift=2mm},
    boxrule=1pt, sharp corners
}

\newtcolorbox{importantbox}[1]{
    enhanced, colback=yellow!10!white, colframe=yellow!80!black, colbacktitle=yellow!80!black,
    coltitle=black, fonttitle=\bfseries, title=#1, sharp corners
}

\newtcolorbox{warningbox}[1]{
    enhanced, colback=red!5!white, colframe=red!75!black, colbacktitle=red!75!black,
    coltitle=white, fonttitle=\bfseries, title=#1, sharp corners
}

\newtcolorbox{dialoguebox}[1]{
    enhanced, breakable, colback=green!4!white, colframe=green!45!black, colbacktitle=green!45!black,
    coltitle=white, fonttitle=\bfseries, title=#1, sharp corners
}

% Listings
\lstset{
    language=Python,
    basicstyle=\ttfamily\small,
    keywordstyle=\color{blue},
    stringstyle=\color{red!60!black},
    commentstyle=\color{green!60!black},
    breaklines=true,
    frame=single,
    numbers=left,
    numberstyle=\tiny\color{gray},
    captionpos=b,
    extendedchars=false
}

% Front-page metadata
\newcommand{\notetitle}{课程笔记：深度研究（Deep Research Agents）}
\newcommand{\noteauthors}{Bowen-mantou（整理）}
\newcommand{\notedate}{2026-10-23}
\newcommand{\videochannel}{B 站「常青藤名校课」· CMU 11-768 AI Agents（英文原声字幕版）}
\newcommand{\videopublishdate}{2026-09-09}
\newcommand{\videoduration}{72:25}
\newcommand{\videourl}{https://www.bilibili.com/video/BV1a6Yx62EH4?p=24}
\newcommand{\videocoverpath}{cover.jpg}

\begin{document}

\begin{titlepage}
\centering
\vspace{1.2cm}
{\huge\bfseries \notetitle\par}
\vspace{0.8cm}
{\large \noteauthors\par}
\vspace{0.3cm}
{\large \notedate\par}
\vspace{1.2cm}

\ifdefempty{\videocoverpath}{
{\small 视频封面因版权未收录，见下方视频链接。\par}
}{
\par
}

\vfill
\begin{tcolorbox}[width=0.9\textwidth, colback=black!2!white, colframe=black!60, sharp corners]
\textbf{视频作者/频道}：\videochannel\par
\textbf{发布日期}：\videopublishdate\par
\textbf{视频时长}：\videoduration\par
\textbf{视频链接}：
\ifdefempty{\videourl}{
未填写
}{
\href{\videourl}{\nolinkurl{\videourl}}
}
\par\textbf{PDF制作Skill}：\href{https://github.com/wdkns/wdkns-skills}{\nolinkurl{https://github.com/wdkns/wdkns-skills}}
\end{tcolorbox}
\end{titlepage}

\tableofcontents
\newpage

%% --- 正文内容开始 --- %%

\section{引言：什么是深度研究 agent}

这一讲由 CMU LTI 的新晋助理教授 Akari Asai 主讲（她称这是自己在 CMU 的第一堂课），主题是 \textbf{deep research agents}——面向复杂信息需求的搜索型 agent。她按「评测（evaluation）→建模（modeling）→检索（retrieval）」三块组织全讲，并在最后讨论开放问题。

\subsection{从「一次检索」到「反复搜索与综合」}

课堂先对比了两个问题。第一个是「我在 CMU 的办公室是几号？」——只要搜一次、找到系里个人主页上的信息就能回答；用一个简单的「检索增强生成（RAG）」pipeline，让模型读一篇文档并抽取答案即可。整个过程是：\textbf{一次搜索、一篇正确网页、直接作答}。

第二个问题是「AI agent 能像人类专家一样综合科学文献吗？」——没有哪一篇文档能直接回答它，必须反复查看多篇论文与多个网站、比较不同研究的结论。此时需要的是 \textbf{deep research agent}：它\textbf{反复规划、调用外部搜索、反思结果}，再决定是否继续搜索，最后综合出带引用的长答案。

\subsection{Deep research agent 的一次迭代}

把上面的过程拆开，agent 在每一轮做四件事：\textbf{计划（plan）}——我需要先找到哪些信息，例如文献综合的研究、语言模型与人类专家对比的研究；\textbf{搜索（search）}——生成搜索 query，交给外部搜索系统，得到一组文档；\textbf{反思（reflect）}——这些结果是否回答了问题？还缺什么？\textbf{修订（revise）}——改写 query 再搜。如此迭代若干轮，收集到足够信息后\textbf{综合出最终答案并附引用}。

这条链路也带来了三类新的研究问题：如何\textbf{评测}这类 agent；如何\textbf{训练}它（从示范到 RL）；以及如何为它构建更好的\textbf{检索系统}——这正是本讲的三大部分。

\subsection{本章小结}
Deep research agent 面向「没有单一文档能回答」的复杂信息需求：它反复「规划—搜索—反思—修订」，最后综合多来源并给出引用。评测、训练与检索是支撑它的三根柱子。

\section{评测：从经典 QA 到 deep research benchmark}

\subsection{经典 QA benchmark 的四个局限}

课堂以一个经典问答 benchmark（Natural Questions，来自真实 Google 搜索 query、用当年维基百科标注答案）为例，指出它与今天 deep research 用法的差距，归纳成四个维度：

\begin{enumerate}
\item \textbf{任务复杂度}：大多是单跳查找（simple lookup），一个事实能在一次搜索里找到；
\item \textbf{领域}：多为通用领域（体育、娱乐），而 deep research 越来越多用于专家领域（金融、医疗、科研）；
\item \textbf{答案质量}：多为短实体答案（short-form），而 deep research 常需要长文开放式综合；
\item \textbf{证据支持}：旧 benchmark 不检查答案是否有可靠证据支撑，而信息类问答恰恰需要知道「答案从哪来」。
\end{enumerate}

\subsection{BrowseComp：多条件、多跳、短答案}

OpenAI 的 \textbf{BrowseComp} 是当前最常用的 deep research benchmark 之一：包含 1000 多道人工撰写的题目。它的构造方式很巧妙——作者\textbf{从已知答案倒推}：先选定一个人物、事件或物品作为目标答案，再收集一组线索（clue），把多条线索组合成一道复杂问题。这样问题的最终答案仍是短实体（便于自动判定是否正确），但解题必须经过多次搜索与多条件筛选。

题目还经过多道验证：先用语言模型加搜索确认「不能只靠模型本身或 5 次 Google 搜索答对」，再让人类确认「10 分钟内答不对」。课堂展示的一道题要求同时满足五条线索（虚构角色、1970 年代电视剧、少于 50 集等），只有同时满足所有条件的才是正确答案。

结果很能说明难度：在带浏览工具的设置下，GPT-4o 只有 0.6\%、GPT-4o+browsing 1.9\%、GPT-4.5 0.9\%、o1（medium）9.9\%；而专门为 deep research 训练的 OpenAI Deep Research 达到 51.5\%，人类参考正确率是 25.3\%。课堂引用原论文的结论：即使有外部搜索，当时的强模型仍大幅落后于专用系统，而人类专家在该 benchmark 上也只有约 25\%。

课堂还引用了 Tongyi DeepResearch 的消融，说明这类任务的\textbf{长度尺度}：把最大上下文长度从 8K 提到 128K，准确率从接近 0 提升到约 40\% 量级；而工具调用上限设到 128 次时效果最好。也就是说，高难度 deep research 需要\textbf{上百次搜索调用与很长的轨迹}。这带来一个很实际的问题：用真实 web 搜索 API 既慢又贵，而且 Google 的结果几乎每天变化，评测难以复现。

\subsection{BrowseComp-Plus：固定语料与可复现评测}

\textbf{BrowseComp-Plus} 解决复现性问题：它\textbf{固定 agent 可检索的文档集合}，把「开放 web 搜索」换成「在受控语料里检索」。构造流程是：从问题与答案出发，用语言模型找出可能有用文档、再由人工核验；然后混入「看起来有用但其实无关」的负样本文档，最终形成一个约 10 万篇文档、约 830 道题的语料库。这样评测可重复，也能单独研究检索系统（这正是第 4 节的内容）。

除了通用域，课堂还介绍了若干专家领域 benchmark：\textbf{FinSearchComp}（金融；由 MBA/金融从业者出题，使用金融数据库并用计算工具核对答案，据课堂口径有 635 道深度研究题、经过多轮专家审核）、\textbf{MedBrowseComp}（医疗；查检验、药物、监管事实）、\textbf{AutoResearchBench}（科研；找出满足多条条件的那篇论文）。共同结论是：高质量的专家域 deep research benchmark 必须靠\textbf{专家出题与核验}。

\subsection{开放式长答案与 rubric 评测}

短答案可以用 exact match 判定，但很多 deep research 任务要求\textbf{长文开放式综合}。传统做法是标注参考并用相似度指标（ROUGE、BERTScore、BLEURT 等）比较模型输出；课堂指出这类指标很不可靠：一项针对长文问答的实证研究让专家在两条回答中选更好的，再与指标比较，结果这些指标与专家偏好的\textbf{一致率只有 58\%（ROUGE）、57\%（BERTScore）、62\%（BLEURT）}，只比随机 50\% 略高。

替代方案是 \textbf{rubric（评分量表）评测}：让专家写一组评价标准，每条标准有具体描述与权重；对每条标准，用一个语言模型判断模型答案是否满足，再对满足项加权求和。课堂（基于其所在团队的 ScholarQA 工作）报告：只要 rubric 质量高，专家之间的一致性约 \textbf{80\%}，语言模型与专家之间约 \textbf{79\%}——这意味着 rubric 可以在保证质量的前提下被语言模型大规模使用。

代价是 rubric 很贵：课堂说一个博士专家标注一条 rubric 大约要 \textbf{1 小时}，难以规模化。于是有工作尝试自动生成 rubric：例如 ResearchQA 从综述论文出发让语言模型生成问题与 rubric，可扩到约 2.1 万道题、16 万条 rubric；但人工审计发现约 \textbf{15\%} 的自动 rubric 存在错误——引用不存在的论文、重复问题、或标准过于空泛以至于无法区分好坏答案。结论是：自动 rubric 必须\textbf{用检索/来源做 grounding}，并经过审计。

此外，\textbf{证据质量}本身也值得评测：DeepResearch Bench 提出 FACT 框架，抽取答案中的「claim + 引用 URL」，逐条检查该 URL 是否真的支持这句话；课堂用一道题的例子说明：两条 claim 引用同一个网页，只有一条被支持，则该题的引用准确率为 50\%。不过课堂也指出：FACT 检查「引用是否支持 claim」，但不检查「被引用的来源本身是否可靠」，后者仍是开放问题。

\subsection{本章小结}
评测经历了从「单跳、短答案、无证据」到「多跳、专家域、长答案、要证据」的升级：BrowseComp 用多条件问题与逆向构造把难度拉满，BrowseComp-Plus 用固定语料换来可复现；开放式答案要靠 rubric（专家—语言模型一致性约 79\%）而不是相似度指标；引用质量可用 FACT 框架逐条核验。课堂的选用建议是：短答案验证型任务用 BrowseComp（质量高但贵）；要降本、要研究检索用 BrowseComp-Plus；开放式任务用 DeepResearch Bench；专家域可从 FinSearchComp 入手。

\section{建模：如何训练 deep research agent}

\subsection{推理链路回顾}

给定一个复杂问题，agent 先产生思考 token 与计划，再生成工具调用（搜索 query），外部搜索系统返回文档并进入上下文；模型基于新证据继续推理，判断是否继续搜索；经过若干轮后产出最终答案与引用。这一节讨论\textbf{怎么训练}这样的模型。

\subsection{三阶段训练配方}

课堂以阿里巴巴通义团队的 DeepResearch 为例，给出常见三阶段配方：

\begin{enumerate}
\item \textbf{Mid-training}：让 base model 在大规模数据上接触 agent 行为，成为「agent-ready」的底座；
\item \textbf{SFT}：让学生模型模仿强教师生成的（成功）轨迹；
\item \textbf{RL}：让策略模型从\textbf{自己的尝试}中学习，用可验证奖励或 rubric 奖励更新。
\end{enumerate}

三个阶段都依赖高质量数据：mid-training 与 SFT 需要高质量\textbf{轨迹}，RL 需要有可验证答案的\textbf{问题}。数据从哪里来？

\subsection{用知识图谱构造多跳问题}

课堂用知识图谱（knowledge graph）解释训练数据的规模化构造：实体是节点、关系是有向边。例如「Arun Newell」与「Carnegie Mellon」两个节点之间有一条「他是该校教授」的边；从这个实体出发，还可以连到「1975 年图灵奖」「与 Herbert Simon 合著 Human Problem Solving」，而 Simon 又连到「诺贝尔经济学奖」。

做法是：在知识图谱上做\textbf{随机游走}，采样一组相连的事实，把它们组合成一道多跳复杂问题，并以「被查询的实体」作为目标答案。这与 BrowseComp 的构造方式同源——某种意义上是在\textbf{为 benchmark 反向工程训练数据}。课堂也诚实指出：这种构造方式训练出的模型能否迁移到其他类型的 deep research benchmark，仍是开放问题；但它提供了一条可扩展的高质量数据构造路径。

有了问题与答案，还需要教师轨迹：让更强的模型（专有模型或更大的开源模型）为每题生成多条轨迹，按「最终答案是否匹配 / 格式是否合法」筛掉失败轨迹，只保留成功轨迹做 SFT。有的系统更早一步做 \textbf{agentic mid-training}：把大量 agent 轨迹混入持续预训练数据，再做同样的下游 SFT；课堂给的对比显示，同样的 SFT 数据下，mid-trained 底座的最终表现明显更好（尤其在复杂 benchmark 上）。

\subsection{短答案任务：RLVR 与 GRPO}

对短答案、可验证的任务，可以直接用 \textbf{RLVR}：策略模型对同一题生成多条轨迹，用 exact match 或语言模型 judge 判断最终答案是否正确，得到奖励，再用策略梯度更新。课堂明确提到这条路线最早在搜索 agent 上的尝试之一就是 Search-R1 一类工作。

以 GRPO 为例：同一道题生成一组轨迹（如 3 条），其中两条答对、一条答错，组均值以上的轨迹被强化、以下的被抑制。课堂用「Herbert Simon（奖励 0）vs Allen Newell（奖励 1，出现两次）」的例子说明组内相对优势：答对的轨迹优势为正、答错的为负，更新时把概率质量推向正确尝试。实际训练还会加 clipping（限制单步更新幅度）与 KL 正则（限制偏离初始策略太远），这两项在大规模训练时对稳定性很重要。

\subsection{工程：让长轨迹 RL 跑得动}

Agentic RL 很贵也不稳定：轨迹长达上百次搜索调用，每次调用要等外部返回。课堂介绍了几条工程手段：\textbf{异步 rollout}（一条轨迹在等工具时，其他轨迹继续生成，完成后统一打分组 batch）、\textbf{缓存搜索结果}、更健壮的 API 客户端与重试、以及用\textbf{本地索引/模拟环境}替代真实 API。训练曲线通常表现为奖励上升而策略熵保持相对稳定——这是训练健康的一个观察点。

\subsection{开放式任务：rubric-based RL 与 DR Tulu}

短答案有 exact match，但开放式研究报告没有二值正确性：一份好的文献综述需要覆盖度、相关性、事实性与恰当引用。课堂介绍了其团队的 \textbf{DR Tulu}（Reinforcement Learning with Evolving Rubrics for Deep Research），核心是\textbf{让 rubric 与策略一起演化}。

流程：先用外部搜索\textbf{grounding} 出一组初始 rubric；训练时对每题采样 8 条轨迹；用语言模型\textbf{对比多条轨迹}，生成新的「正 rubric」（奖励好的行为）与「负 rubric」（惩罚具体错误）；维护一个 rubric 缓冲区，用它给最好的回答打分并更新策略。关键在于：反馈不是训练开始时固定的，而是\textbf{随策略探索不断捕捉新的区分}——策略变强后，rubric 也能识别更细的质量差异。

课堂报告的消融很值得记住：用\textbf{随机奖励}做 online RL 看不到任何提升；用\textbf{初始 search-grounded rubric} 明显更好；而 evolving rubric 在训练更久时收益更大——说明「初始 rubric 质量」与「rubric 的适应性」两者都重要。另一个关键发现是 \textbf{SFT 是 RL 的热身}：从 base 模型直接 RL（600 步）只有约 40\%，而只用约 5\% 的 SFT 数据做「欠训练 SFT」后，同样步数可再提升约 10 个百分点；完整 SFT 则给出更强的 RL 起点。

最终，DR Tulu-8B 在四个 deep research benchmark 的平均分上，以远低于专有系统的推理成本取得有竞争力的表现（课堂的成本—性能图：横轴是每查询成本，纵轴是平均分；DR Tulu-8B 位于左上角）。

\begin{warningbox}{rubric 生成不能裸用语言模型}
课堂引用 ResearchQA 的审计：纯语言模型生成的 rubric 约 15\% 有错（引用不存在的论文、重复问题、标准空泛）。可用的 rubric 需要外部证据 grounding + 审计；rubric 生成模型越强越好（课堂口径：GPT-4.1 与 Qwen3 系做 rubric 生成约差 10 个点），开源 rubric 模型仍是开放问题。
\end{warningbox}

\subsection{本章小结}
训练配方是 mid-training → SFT → RL：用知识图谱随机游走构造多跳问答、用强教师轨迹做 SFT、用 RLVR/GRPO 从自身尝试学习。开放式报告用 rubric-based RL，DR Tulu 让 rubric 与策略共同演化；SFT 是 RL 的热身，少量 SFT 也能显著抬高 RL 起点；rubric 必须 grounded 且审计。

\section{检索：为 agent 构建搜索基础}

\subsection{本地检索与黑盒 API}

课堂把 deep research agent 可用的检索分成两类：

\begin{itemize}
\item \textbf{本地检索（local search）}：已知语料时，可在本地建词法检索（BM25）或神经 embedding 检索（训练过的相似度模型）。例如 Search-R1 一类工作用 2018 年维基百科语料 + embedding 检索器。
\item \textbf{黑盒 API}：开放 web 场景用 web search API（如 Serper）或浏览 API；科研场景用 Semantic Scholar、Google Scholar 等论文搜索 API。
\end{itemize}

选择取决于任务：要做好 BrowseComp/公开 web 题，需要真实 web 搜索；要做 BrowseComp-Plus 这类受控语料评测，本地索引更可复现也更便宜。

\subsection{Embedding 检索的复习与训练目标}

标准 embedding 检索分两步：离线用文档编码器把所有文档编码成向量、建向量索引；在线用查询编码器编码 query，做近邻检索。训练常用\textbf{对比学习}：每个 query 配一个正文档与若干负文档，目标函数提高正文档相对负文档的得分：

$$
\mathcal{L} = -\log \frac{\exp\big(s(q, d^+)\big)}{\exp\big(s(q, d^+)\big) + \sum_{d^-} \exp\big(s(q, d^-)\big)},
$$

其中 $q$ 是查询，$d^+$ 是正文档，$d^-$ 是负文档，$s(\cdot,\cdot)$ 是相似度打分（如内积或余弦）。负样本可以来自 \textbf{in-batch negatives}（同一 batch 里其他 query 的正文档充当负例，DPR 的经典做法），也可以由较弱的检索器（如 BM25）挖出\textbf{困难负例}再混合。

\subsection{Agent 检索的特殊性：检索器只看得到最后一个 query}

标准检索器是为「人类用户的一次查询」设计的；但 deep research agent 在每次检索时拥有\textbf{很丰富的历史}：原始任务、之前的推理 token、之前的 query、之前搜索的结果。课堂指出一个关键问题：现有系统通常把\textbf{当前 query}单独交给检索器，历史信息被丢掉了。

课堂介绍了其团队的 \textbf{AgentIR}：把「当前推理 token + query」一起输入检索器，并针对这种「推理感知（reasoning-aware）检索」训练 embedding。消融显示两点都重要——只在推理时拼接推理 token 有小幅提升，而对检索器做针对性训练能带来更大提升；两者结合效果最好。另一个有意思的发现是\textbf{上下文不是越多越好}：把整条轨迹都塞给检索器反而变差（上下文太长），而只用「上一步的推理 token」最有效——它恰好总结了任务与上一步搜索结果。只用历史 query（不含搜索结果）帮助有限。

课堂报告：AgentIR 在 BrowseComp-Plus 上同时提升准确率、降低平均搜索调用次数；它只有 4B，比 Qwen3-8B 这类通用 embedding 模型更小；训练数据只来自 BrowseComp-Plus 的少量轨迹，却能在不同 agent 模型与不同 benchmark 上泛化——说明「为 agent 的检索需求专门训练检索器」是一条值得投入的方向。

\subsection{本章小结}
检索要么在本地受控语料上建 BM25/embedding 索引（可复现、便宜），要么调黑盒 web/论文搜索 API（开放但慢且不可复现）。embedding 用对比学习训练，负样本可用 in-batch 与 BM25 困难负例。Agent 检索的特殊性在于历史丰富而检索器只看最后一个 query；AgentIR 把「当前推理 + query」交给检索器并专门训练，在更少搜索次数下取得更高准确率，且小模型也能泛化。

\section{总结与延伸}

主讲人最后用四点总结全讲：(1) \textbf{任务}——deep research 是「迭代搜索 + 证据综合」，不是单次检索问答；(2) \textbf{评测}——要同时评搜索复杂度、答案质量、引用支撑，以及 rubric 本身的质量；(3) \textbf{训练}——用可验证问题构造任务与轨迹，从示范学习（mid-training/SFT），再从自身尝试中学习（RL）；(4) \textbf{检索}——要为 agent 当前的检索需求构建并训练检索器，而不是沿用为人类查询设计的检索器。

全讲可以压缩为六条：

\begin{enumerate}
\item Deep research agent 的循环是「规划 $\rightarrow$ 搜索 $\rightarrow$ 反思 $\rightarrow$ 修订 $\rightarrow$ 综合带引用答案」，适用于没有单一文档能回答的复杂信息需求。
\item 经典 QA benchmark 在四个维度上不够用（复杂度、领域、答案形态、证据支持）；BrowseComp 用逆构造的多条件问题把难度拉满，BrowseComp-Plus 用固定语料解决可复现性，FinSearchComp 等专家域 benchmark 依赖专家出题核验。
\item 开放式答案不能用相似度指标评（与专家偏好一致率仅 57\%--62\%），rubric 评测在高质量前提下专家—模型一致性约 79\%；自动 rubric 有约 15\% 错误，需要 grounding 与审计；引用质量可用 FACT 逐条核验。
\item 训练配方是 mid-training $\rightarrow$ SFT $\rightarrow$ RL：知识图谱随机游走构造多跳问题、强教师轨迹筛选后 SFT、RLVR/GRPO 从自身尝试学习；SFT 是 RL 的热身（少量 SFT 也能显著抬高起点）。
\item 开放式报告用 rubric-based RL，DR Tulu 让 rubric 与策略共同演化，用 8B 模型在低成本下取得有竞争力的平均分。
\item 检索是 deep research 的底座：本地 BM25/embedding 与黑盒 API 各有取舍；AgentIR 表明把「当前推理 + query」交给检索器并专门训练，可以在更少搜索次数下提高准确率。
\end{enumerate}

\begin{importantbox}{一句话记住本讲}
\textbf{Deep research = 评测（多跳/专家域/长答案/证据）、训练（mid-training $\rightarrow$ SFT $\rightarrow$ 可验证或 rubric 的 RL）、检索（为 agent 的推理历史专门训练）三件事的组合；三者缺一，agent 都会在「搜不到、评不准、学不会」上卡住。}
\end{importantbox}

\textbf{延伸思考}：(1) 从系统角度看，deep research 的成本主要由「搜索调用次数」与「上下文长度」决定，而两者又由训练配方决定——这与第 9 讲「RL 的组大小、全对全错比例」一样，属于「采样预算设计」问题；(2) rubric 与第 8 讲 SFT 的「教师信号」、第 9 讲的「奖励设计」是同一条主线：开放式任务缺少可验证信号时，只能用「可审计的近似信号」，而近似信号的偏差会被策略利用（reward hacking 的另一种形态）；(3) retrieval 与 agent 的联合训练很可能是下一个热点——检索器需要理解 agent 的意图与历史，而 agent 又需要理解检索器的偏好，两者目前还是分开训练的。课堂在问答中还给了实操建议：从 8B 级模型（Qwen3-8B 一类）与稳定 web search API 起步，训练时用 BrowseComp-Plus 的本地索引（BM25 或 Qwen3-4B/8B embedding）压成本；embedding 选型可参考 MTEB；报告结果时最好在多个 agent 模型（如 GPT-OSS、GLM）上验证，而不是只调优单一 agent。

\end{document}
